\section{Degradación de autoridad: el certificado como regularización de un problema singular}
\label{sec:autoridad}

La \cref{obs:autoridad} señaló que un par sin autoridad instantánea no puede
corregir su barrera, pero lo hizo de forma binaria: el coeficiente de control se
anula o no se anula. La situación real es continua, y su descripción exacta
cambia el criterio de diseño.

\subsection{Los tres regímenes}

Sea $A=Q^{1/2}G_k$ con descomposición en valores singulares $A=U\Sigma V^{\!\top}$
y $\tilde e=Q^{1/2}\Wdem_k$. Con el cambio ortogonal $\bm\lambda=V\bm\mu$, el
programa \eqref{eq:qp} se desacopla en problemas escalares independientes, uno
por dirección singular:
\begin{equation}
  \min_{\mu_j}\ \tfrac12\bigl(\sigma_j\mu_j-\tilde e_j\bigr)^{2}
  +\tfrac{\varepsilon_\lambda}{2}\,\mu_j^{2}.
  \label{eq:escalar}
\end{equation}

\begin{teorema}[Degradación de autoridad]
\label{th:autoridad}
En el régimen interior de \eqref{eq:qp}, para cada dirección singular $j$:
\[
  \mu_j^{\star}=\frac{\sigma_j\,\tilde e_j}{\sigma_j^{2}+\varepsilon_\lambda},
  \qquad
  \text{residuo}_j=\sigma_j\mu_j^{\star}-\tilde e_j
  =-\frac{\varepsilon_\lambda\,\tilde e_j}{\sigma_j^{2}+\varepsilon_\lambda}.
\]
En consecuencia:
\begin{enumerate}[label=(\roman*),leftmargin=2.2em,topsep=0.2em,itemsep=0.15em]
  \item \emph{autoridad plena} si $\sigma_j^{2}\gg\varepsilon_\lambda$:
  $\mu_j^{\star}\to\tilde e_j/\sigma_j$ y el residuo tiende a cero, con esfuerzo
  que crece como $1/\sigma_j$;
  \item \emph{parálisis óptima} si $\sigma_j^{2}\ll\varepsilon_\lambda$:
  $\mu_j^{\star}\to\sigma_j\tilde e_j/\varepsilon_\lambda\to0$ y el residuo
  tiende a $-\tilde e_j$. El programa es estrictamente convexo y su solución es
  única y óptima, y sin embargo \emph{renuncia} a corregir el error;
  \item el umbral está en $\sigma^{\dagger}=\sqrt{\varepsilon_\lambda}$, donde el
  residuo es exactamente la mitad de la demanda proyectada.
\end{enumerate}
\end{teorema}

\begin{proof}
Derivando \eqref{eq:escalar} respecto a $\mu_j$ e igualando a cero se obtiene
$\sigma_j(\sigma_j\mu_j-\tilde e_j)+\varepsilon_\lambda\mu_j=0$, de donde
$\mu_j^{\star}(\sigma_j^{2}+\varepsilon_\lambda)=\sigma_j\tilde e_j$. La
expresión del residuo resulta de sustituir. Los tres regímenes son los límites
$\sigma_j^2/\varepsilon_\lambda\to\infty$, $\to0$ y $=1$.
\end{proof}

\begin{observacion}[La parálisis óptima no es un fallo del solver]
\label{obs:paralisis}
El caso (ii) merece énfasis porque contradice la intuición: un QP correctamente
planteado, convexo y resuelto con exactitud, decide no mover la carga. No hay
error numérico ni infactibilidad. La regularización que garantiza unicidad
—necesaria para el \cref{co:agarre} y para la contracción del solver— es la
misma que produce la renuncia cuando la autoridad es escasa. Es un compromiso
estructural del planteamiento, no un defecto de implementación.
\end{observacion}

\begin{observacion}[La regularización no solo estabiliza: elige]
\label{obs:tikhonov-selecciona}
Conviene añadir algo que el enunciado anterior deja implícito y que mejora la
lectura del parámetro $\varepsilon_\lambda$. La regularización empleada es de
tipo Tikhonov, y ese esquema tiene una propiedad de \emph{selección}: cuando el
problema sin regularizar admite infinitas soluciones, la solución regularizada
converge, al hacer $\varepsilon_\lambda\to0$, a la de \emph{norma mínima} entre
todas las óptimas \cite{clempner2024optimization}.

Comprobación sobre un sistema de rango deficiente con infinitas soluciones: la
distancia a la solución de norma mínima es $1{,}1\times10^{-1}$ para
$\varepsilon=10^{-1}$, $1{,}2\times10^{-3}$ para $10^{-3}$ y
$1{,}4\times10^{-9}$ para $10^{-9}$.

La lectura para este trabajo es que $\varepsilon_\lambda$ cumple tres funciones a
la vez, no una: garantiza unicidad, produce la parálisis de
\cref{th:autoridad}(ii) cuando la autoridad escasea, y —cuando la autoridad
sobra y el reparto es indeterminado— selecciona el de \emph{menor esfuerzo total}.
Esa tercera función es deseable y merece enunciarse: en una coalición
sobreactuada, donde muchos repartos realizan la misma demanda, el mecanismo no
elige uno arbitrario sino el más barato.
\end{observacion}

\begin{teorema}[Valor exacto de un actuador adicional]
\label{th:actuador-adicional}
Sea $V(K)=\tfrac12 d^{\!\top}K^{-1}d$ el coste óptimo del reparto para una
demanda $d$, y sea un candidato opcional que añade un canal escalar $b$ con coste
propio $c\,u^{2}/2$, $c>0$, de modo que $K^{+}=K+bb^{\!\top}/c$. La disminución
óptima de coste es exactamente \cite{mayorga2026m1}
\[
  V(K)-V(K^{+})\;=\;\frac{\bigl(b^{\!\top}K^{-1}d\bigr)^{2}}
  {2\,\bigl(c+b^{\!\top}K^{-1}b\bigr)}\ \ge\ 0 .
\]
\end{teorema}

\begin{proof}
Por Sherman--Morrison,
$(K+bb^{\!\top}/c)^{-1}=K^{-1}-K^{-1}bb^{\!\top}K^{-1}/(c+b^{\!\top}K^{-1}b)$;
sustituyendo en $V$ se obtiene la expresión. El denominador es positivo porque
$c>0$ y $K\succ0$.
\end{proof}

\begin{observacion}[Qué hace valioso a un recluta]
\label{obs:actuador-valor}
La fórmula, verificada a precisión de máquina sobre instancias aleatorias, no es
una heurística y separa dos cosas que se confunden. El numerador
$(b^{\!\top}K^{-1}d)^{2}$ mide la \emph{alineación} del canal con la dirección que
la demanda encarece; el denominador incorpora el coste propio y la
\emph{redundancia} con lo ya disponible, $b^{\!\top}K^{-1}b$. Un canal de gran
capacidad nominal en una dirección que no importa, o que duplica lo que la
coalición ya tiene, vale poco; uno modesto que corrige la dirección cara vale
mucho.

Es la versión cuantitativa de \cref{obs:brazo} —donde toca importa tanto como
cuánto empuja— y da al reclutamiento un criterio exacto de admisión marginal:
el valor de un candidato es un número calculable a partir de $K^{-1}d$, que la
coalición ya posee. También precisa la lectura de \cref{obs:shapley}: Shapley
premia ser pivote en muchas coaliciones; este teorema premia ser útil en
\emph{esta}, y ambos criterios pueden discrepar sobre el mismo AMR.
\end{observacion}

\begin{observacion}[La misma selección sesga hacia la inacción]
\label{obs:tikhonov-sesgo}
Conviene mirar esa propiedad con desconfianza allí donde el objeto seleccionado
no es un esfuerzo sino una decisión. La regularización elige por \emph{norma
mínima}, no por bienestar, y en un problema de formación de coalición la
alternativa de norma mínima es «nadie se une».

La comprobación matiza el riesgo y conviene reproducirla con precisión. Sobre un
potencial de doble pozo con el equilibrio activo estrictamente mejor, la
regularización lo selecciona correctamente para $\varepsilon\le0{,}5$, y solo con
$\varepsilon=2$ se pasa al pasivo. Pero si los dos pozos están igualados, la
regularización rompe el empate a favor del pasivo para \emph{todo}
$\varepsilon>0$ ensayado, de $0{,}5$ a $10^{-3}$.

Hay por tanto una ventana: $\varepsilon$ suficientemente pequeño para no invertir
la preferencia, y suficientemente grande para dar el condicionamiento que se
buscaba. Y hay un patrón común con \cref{th:autoridad}: la
regularización que produce la parálisis óptima de \cref{th:autoridad}(ii) en la
capa de esfuerzo produciría, aplicada a la capa de decisión, la misma inclinación
a no hacer nada. Ambas son la misma preferencia por la norma pequeña, y en las
dos hay que decidir a sabiendas.
\end{observacion}

\subsection{El caso singular}

Cuando $\sigma_j=0$ con $\tilde e_j\neq0$, el término de tarea deja de depender
del control en esa dirección: no existe solución clásica que anule el residuo.
La familia $\{\bm\mu^{\star}(\varepsilon_\lambda)\}_{\varepsilon_\lambda>0}$ es
entonces una \emph{sucesión minimizante} del problema singular, en el sentido
preciso de \cite{glizer2022singular}, donde la solución de un problema
extremal singular no es un elemento sino una sucesión, y la regularización de
control barato es el método constructivo para obtenerla.

El certificado de \eqref{eq:qp} es, por tanto, la regularización de un problema
que en la frontera es singular. Esa lectura tiene una consecuencia práctica:
cuando $\sigma_j\to0$, no procede refinar el solver ni bajar $\varepsilon_\lambda$
esperando convergencia, porque el límite no existe como control admisible. Lo
que procede es rechazar la configuración.

\begin{corolario}[La guarda de autoridad se acopla a la regularización]
\label{co:guarda}
Un criterio de admisión de la forma $\sigma_{\min}(Q^{1/2}G_k)\ge\sigma_{\mathrm{tol}}$
con $\sigma_{\mathrm{tol}}$ elegido de forma independiente es incoherente. La
guarda debe exigir
\[
  \sigma_{\min}\bigl(Q^{1/2}G_k\bigr)\ \ge\ \kappa\sqrt{\varepsilon_\lambda},
  \qquad \kappa>1,
\]
de modo que el residuo relativo quede acotado por $1/(\kappa^{2}+1)$. La
profundidad de la guarda y la regularización del programa son un solo parámetro
de diseño, no dos.
\end{corolario}

\begin{proof}
Sustituyendo $\sigma_j=\kappa\sqrt{\varepsilon_\lambda}$ en la expresión del
residuo de \cref{th:autoridad} se obtiene
$|\text{residuo}_j|/|\tilde e_j|=\varepsilon_\lambda/(\kappa^{2}\varepsilon_\lambda+\varepsilon_\lambda)
=1/(\kappa^{2}+1)$, y la cota es decreciente en $\sigma_j$.
\end{proof}

\begin{observacion}[Alcance]
\label{obs:autoridad-alcance}
\Cref{th:autoridad} describe el régimen interior, en el que las restricciones de
cono y de caja no están activas; con restricciones activas la solución queda
caracterizada por las condiciones KKT y el conjunto activo. La lectura cualitativa —degradación continua con umbral en
$\sqrt{\varepsilon_\lambda}$— se conserva, pero las expresiones cerradas no.
Tampoco se afirma nada sobre el juego diferencial completo: el resultado es del
problema congelado, que es donde \cite{glizer2022singular} delimitan su
propio análisis.
\end{observacion}
